\begin{abstract}
Benchmark racing---the standard method for choosing among LLM agent designs---conflates per-step cost, model competence, information structure, and routing quality into a single opaque comparison.
We show that, under a packed latent-mode family assumption, the performance gap between any two agent designs decomposes exactly into four independently measurable terms.
The decomposition is an algebraic identity, not an approximation.

The decomposition yields a closed-form crossover formula---a small model with $D$ retries matches a large model when $\Dcross = \log(1-p_l)/\log(1-p_s)$, and dominates when this depth falls below the cost ratio---and a monotonicity theorem: stronger models need less scaffolding.
In synthetic evaluations, decomposition-guided search uses $50\times$ fewer evaluations than benchmark racing at comparable regret.
A Bernstein corollary makes the required pilot estimation practical for hard tasks.

An honest negative result clarifies scope: the four terms are diagnostic performance accounting, not Lagrangian shadow prices.
The decomposition tells a builder which lever matters most but does not solve for the optimum in closed form.
The framework applies under a packed-family assumption on the task distribution; its scope and limitations are characterized explicitly.
\end{abstract}

\begin{keywords}
agent design, certified time, performance decomposition, test-time compute, model selection, routing
\end{keywords}
